On dissout dans l'eau~:
\begin{itemize}
    \item $\qty{117}{\mg}$ de chlorure de sodium
          $(M_{\ce{NaCl}}=\qty{58.5}{\g\per\mol})$~;
    \item $\qty{119}{\mg}$ de bromure de potassium
          $(M_{\ce{KBr}}=\qty{119}{\g\per\mol})$~;
\end{itemize}
pour obtenir une solution de volume $V=\qty{2.0}{\L}$.
\begin{enumerate}
    \item Calculer la concentration molaire des ions dans la solution.
    \item En utilisant les données du tableau page 59 du cours, calculer la
          conductivité de la solution.
    \item Que vaudrait la conductivité si la solution était diluée au dixième~?
\end{enumerate}


